\section{Conclusion}\label{sec:conclusion}

% W = What proposed
This paper presented FMPA, a domain AI agent that helps consumers judge whether
a vegetable price they see today is likely to change in the next 7 or 14 days.
It addresses three challenges: public averages are not store prices, retail
prices are sticky with rare jumps, and LLMs should not invent numbers.

% X = eXcel (where it fits)
FMPA is designed for the moment of purchase: it anchors forecasts on the
observed store price, uses fixed calendar horizons that can later be checked,
and keeps every number inside deterministic MCP tools, with explicit errors
when data are missing.

% Y = Yields
For three vegetables in 13 public markets in Taichung City, RandomForest and
HistGradientBoosting models reached direction macro-F1 of 0.422 and 0.417,
above the no-change baseline but below momentum (0.423 and 0.431), and their
relative-change MAE was worse than that of the no-change baseline. The current
value of FMPA therefore lies in its transparent, store-anchored presentation of
forecasts and baselines, not in superior accuracy.

% Z = Zen (limitations and future work)
The main limitations are that public market prices are only a proxy for
individual store prices and that the benefit of mode B is still unproven.
Future work will collect matured same-store prices for paired validation of
modes A and B; add feature engineering and event
indicators, such as typhoon warnings and arrival volumes; attach uncertainty
estimates to every forecast; and run a user study of the agent interface.
